\documentclass[tikz,border=8pt]{standalone}
\usepackage{tikz}
\usepackage{amsmath}
\usepackage{amssymb}
\usepackage{pgfplots}
\pgfplotsset{compat=1.18}
\usepackage{ctex}
\usetikzlibrary{calc,positioning,arrows.meta,trees}

% LC 236 最近公共祖先（Lowest Common Ancestor）
% 树：[3,5,1,6,2,0,8,null,null,7,4]
% p=5（红色高亮），q=1（红色高亮），LCA=3（绿色高亮）

\begin{document}
\begin{tikzpicture}[
  every node/.style={},
  level distance=1.5cm,
  level 1/.style={sibling distance=5.0cm},
  level 2/.style={sibling distance=2.4cm},
  level 3/.style={sibling distance=1.3cm},
  edge from parent/.style={draw,->,>=Stealth,thick}
]

%% ── 标题 ──────────────────────────────────────────────────
\node[font=\large\bfseries] at (0, 2.2) {LC 236 最近公共祖先（LCA）};
\node[font=\small,gray] at (0, 1.65)
  {目标：$p=5$（红）$,\ q=1$（红），LCA$\ =3$（绿框）};

%% ── 主树 ─────────────────────────────────────────────────
% 根 3 普通节点，LCA 用绿色外框圆圈标注而非 fill
\node[draw,circle,minimum size=0.82cm,fill=blue!10,
      font=\bfseries] (root) at (0,0) {3}
  child {
    % 节点 5：p，红色外框
    node[draw,circle,minimum size=0.82cm,fill=red!8,
         font=\bfseries] (n5) {5}
    child {
      node[draw,circle,minimum size=0.78cm,fill=blue!10,font=\bfseries] (n6) {6}
    }
    child {
      node[draw,circle,minimum size=0.78cm,fill=blue!10,font=\bfseries] (n2) {2}
      child {
        node[draw,circle,minimum size=0.78cm,fill=blue!10,font=\bfseries] (n7) {7}
      }
      child {
        node[draw,circle,minimum size=0.78cm,fill=blue!10,font=\bfseries] (n4) {4}
      }
    }
  }
  child {
    % 节点 1：q，红色外框
    node[draw,circle,minimum size=0.82cm,fill=red!8,
         font=\bfseries] (n1) {1}
    child {
      node[draw,circle,minimum size=0.78cm,fill=blue!10,font=\bfseries] (n0) {0}
    }
    child {
      node[draw,circle,minimum size=0.78cm,fill=blue!10,font=\bfseries] (n8) {8}
    }
  };

%% ── 目标节点额外高亮圆环 ─────────────────────────────────
\draw[red!70!black,very thick] (root) circle (0.44cm);  % LCA=3 绿环
\draw[green!60!black,very thick,dashed] (root) circle (0.50cm);
\draw[red!70!black,very thick] (n5) circle (0.44cm);   % p=5 红环
\draw[red!70!black,very thick] (n1) circle (0.44cm);   % q=1 红环

%% ── 返回值标注（固定绝对坐标，紧靠节点） ────────────────
% 节点 5 左侧：找到 p，向上返回 node5
\node[draw=red!50,fill=red!5,rounded corners=2pt,inner sep=3pt,
      font=\scriptsize,text=red!80!black,align=center]
  at (-6.5,-1.5) {找到 \textbf{p=5}\\返回 \textbf{node5}};
\draw[->,>=Stealth,red!50,thin] (-5.4,-1.5) -- (-2.9,-1.5);

% 节点 1 右侧：找到 q，向上返回 node1
\node[draw=red!50,fill=red!5,rounded corners=2pt,inner sep=3pt,
      font=\scriptsize,text=red!80!black,align=center]
  at (6.5,-1.5) {找到 \textbf{q=1}\\返回 \textbf{node1}};
\draw[->,>=Stealth,red!50,thin] (5.4,-1.5) -- (2.9,-1.5);

% 叶子返回 null（小标签）
\node[font=\tiny,gray] at (-5.0,-3.1) {\texttt{null}};
\node[font=\tiny,gray] at (-2.8,-3.1) {\texttt{null}};
\node[font=\tiny,gray] at (-1.3,-4.6) {\texttt{null}};
\node[font=\tiny,gray] at (0.1,-4.6)  {\texttt{null}};
\node[font=\tiny,gray] at (1.8,-3.1)  {\texttt{null}};
\node[font=\tiny,gray] at (3.8,-3.1)  {\texttt{null}};

% 根 3 顶部：LCA 结论
\node[draw=green!60!black,fill=green!5,rounded corners=3pt,inner sep=4pt,
      font=\small,text=green!60!black,align=center]
  at (0, 0.9) {左=node5，右=node1\\两者非 null $\Rightarrow$ \textbf{LCA = 3}};

%% ── 递归路径虚线箭头（从子到根） ────────────────────────
\draw[->,>=Stealth,red!60,dashed,thick]
  (-2.4,-1.0) -- (-0.5,-0.1)
  node[midway,above left,font=\scriptsize,red!70!black] {5};

\draw[->,>=Stealth,red!60,dashed,thick]
  (2.4,-1.0) -- (0.5,-0.1)
  node[midway,above right,font=\scriptsize,red!70!black] {1};

%% ── 算法说明 ─────────────────────────────────────────────
\node[draw=gray!50,fill=white,thin,rounded corners=3pt,
      font=\small,inner sep=6pt,align=left]
  at (0,-6.5)
  {\textbf{递归规则}：\\
   \quad 1. \texttt{root == null}：返回 \texttt{null}\\
   \quad 2. \texttt{root == p} 或 \texttt{root == q}：返回 \texttt{root}（找到目标）\\
   \quad 3. 递归左子树 $\to$ \texttt{left}；递归右子树 $\to$ \texttt{right}\\
   \quad 4. \texttt{left != null \&\& right != null}：\texttt{root} 即为 LCA\\
   \quad 5. 否则返回非 null 的一侧（目标在同侧子树中）};

\end{tikzpicture}
\end{document}
